\section{Limitations and Threats to Validity}
\label{sec:limitations}

\paragraph{Statistical scope.}
The evaluation has 800 queries but only eight top-level map clusters.  The
hierarchical bootstrap respects this structure, yet eight maps cannot support
precise population claims.  Two maps per family show useful within-family
replication, not coverage of all instances in that family.  The 0.053
bootstrap fraction is descriptive, not a $p$-value, and the interval is not
evidence of equivalence.

\paragraph{External validity.}
All problems are undirected, unit-cost, four-neighbor grids from four known
families.  Results may change on weighted or directed graphs, road networks,
diagonal movement, dynamic tasks, or other pivot-distance representations.
Families occur on both sides of the map-disjoint split, so this tests new maps
within known families rather than unseen-domain generalization.

\paragraph{Fixed design choices.}
$K=4$, $M=32$, and deterministic farthest-first placement form one prespecified
configuration.  The prefix may be unusually informative in some topologies
and redundant in others.  We did not compare alternate prefix sizes, pivot
orders, compression, adaptive stopping, or active-landmark techniques.  This
restriction prevents tuning to the evaluation maps but limits the conclusions to
the specified schedule.

\paragraph{Implementation and measurement.}
Packed-array reads and priority-queue operations have Python- and
hardware-specific costs.  Short searches remain sensitive to OS scheduling,
cache state, allocator behavior, and timer resolution despite a warm-up,
balanced order positions, and eight-run medians.  Disabling stage timers
avoids one perturbation but prevents assigning nanoseconds cleanly to prefix,
suffix, and heap operations.  Exact counters are more portable than runtime
crossover points.

\paragraph{Internal validity.}
A shared engine and trace digest greatly reduce implementation confounds, but
they also share any common bug.  Independent BFS, legal-path validation,
exhaustive small-map tests, immutable source hashes, deterministic replays,
and strict artifact loaders provide independent layers.  Expansion-trace
equality is an implementation invariant under the specified tie breaker; it is
not a general theorem that every Lazy \Astar{} implementation must duplicate
eager order.

\paragraph{Novelty.}
Every algorithmic ingredient has close prior art.  ALT and differential
heuristics establish landmark bounds; subset work establishes that a fixed
selection can trade evaluation cost and informativeness; Lazy \Astar{}
establishes delayed evaluation and notes straightforward multi-stage
extension.  This experiment studies their intersection.  A targeted search
found no exact match, although differently named or unpublished work may
exist.

\paragraph{Analysis rerun.}
The first complete development/evaluation attempt is preserved separately.
Before statistical output was produced, analysis stopped on a representation
mismatch: the loaded schedule used tuples while the equivalent JSON plan used
lists.  We normalized the container types, added a regression test, and reran
the full experiment because the analysis source hash had been fixed before the
first run.  The algorithm, protocol, queries, $K/M$ values, and estimands were
unchanged.  The 800 deterministic summaries match the first attempt; only
timings were remeasured.  The correction was not prompted by the observed
performance.
